\documentclass{article}

\usepackage{float}
\restylefloat{table}

\usepackage{booktabs}

\title{Team Productivity: POC\\\progname}

\author{\authname}

\date{}

\input{../Comments.text}
\input{../Common.text}

\begin{document}

\maketitle

This document summarizes the contributions of each team member up to the POC
Demo. The time period of interest is the time between the beginning of the term
and the POC demo.

\section{Demo Plans}

The two feasibility risks involving \progname{} are the following:

\begin{itemize}
    \item\textbf{Risk 1: Model cannot be split.} The core premise of \progname{} is that it will split up a large language model across more than one device. This will be done by splitting the layers of the model.
    If this cannot be completed, the system will fail.
    \item\textbf{Risk 2: Model can be split, yet provides no improvement.} In the case a model can be split across more than one device, the new device network must improve upon the capabilities that only one device may have.
    The metric the team will initially target is capacity. If two devices can carry a larger model (more parameters) than one of the devices, then there is benefit to using Lake because a model with more parameters is
    generally more performant. If this cannot be done, then users have no reason to begin using \progname{}.
\end{itemize}

In order to address these risks, the team will demonstrate a large language model that is split up among two devices. This model will be larger than either individual device can hold in its memory.
The demonstration session will consist of the team partitioning the model across the two devices using the \progname{} system, followed by starting an inference session by prompting the model.
This will show that both key risks have been mitigated.

\section{Team Meeting Attendance}

\wss{For each team member how many team meetings have they attended over the
time period of interest.  This number should be determined from the meeting
issues in the team's repo.  The first entry in the table should be the total
number of team meetings held by the team.}

\begin{table}[H]
\centering
\begin{tabular}{ll}
\toprule
\textbf{Student} & \textbf{Meetings}\\
\midrule
Total & Num\\
Ethan Walsh & Num\\
Colin Chambachan & Num\\
Matthew Nesbitt & Num\\
Jinwoo Hong & Num\\
Gary Qin & Num\\
\bottomrule
\end{tabular}
\end{table}

\wss{If needed, an explanation for the counts can be provided here.}

\section{Supervisor/Stakeholder Meeting Attendance}

The following section is not applicable to the \progname{} project since it does not have a supervisor.

\begin{table}[H]
\centering
\begin{tabular}{ll}
\toprule
\textbf{Student} & \textbf{Meetings}\\
\midrule
Total & 0\\
Ethan Walsh & 0\\
Colin Chambachan & 0\\
Matthew Nesbitt & 0\\
Jinwoo Hong & 0\\
Gary Qin & 0\\
\bottomrule
\end{tabular}
\end{table}

\section{Lecture Attendance}

\begin{table}[H]
\centering
\begin{tabular}{ll}
\toprule
\textbf{Student} & \textbf{Lectures}\\
\midrule
Total & 11\\
Ethan Walsh & 11\\
Colin Chambachan & 7\\
Matthew Nesbitt & 11\\
Jinwoo Hong & 7\\
Gary Qin & 10\\
\bottomrule
\end{tabular}
\end{table}

\section{TA Document Discussion Attendance}

\noindent \textbf{TA's Name: Rebecca Di Filippo}

\begin{table}[H]
\centering
\begin{tabular}{ll}
\toprule
\textbf{Student} & \textbf{Meetings}\\
\midrule
Total & 4\\
Ethan Walsh & 4\\
Colin Chambachan & 4\\
Matthew Nesbitt & 4\\
Jinwoo Hong & 4\\
Gary Qin & 4\\
\bottomrule
\end{tabular}
\end{table}

\section{Commits, Lines Added and Lines Deleted, Pull Requests}

\begin{table}[H]
\centering
\begin{tabular}{lllllll}
\toprule
\textbf{Student} & \textbf{Commits} & \textbf{Percent} & \textbf{Lines Added} &
\textbf{Lines Deleted} & \textbf{Pull Requests}\\
\midrule
Total & 212 & 100\% & 49811 & 6091 & 74\\
Ethan Walsh & 31 & 14.6\% & 1355 & 505 & 10\\
Colin Chambachan & 83 & 39.2\% & 6103 & 2564 & 28\\
Matthew Nesbitt & 38 & 17.9\% & 13732 & 350 & 18\\
Jinwoo Hong & 25 & 11.8\% & 446 & 254 & 7\\
Gary Qin & 35 & 16.5\% & 2837 & 2418 & 11\\
\bottomrule
\end{tabular}
\end{table}

The reason that Colin and Matthew have significantly higher additions than other team members is due to their addition of some of the initial infrastructure of the codebase. For example,
Matthew added the CI infrastructure and .yaml files.

\section{Issue Tracker}

\wss{For each team member how many issues have they authored (including open and
closed issues (O+C)) and how many have they been assigned (only counting closed
issues (C only)) over the time period of interest.}

\begin{table}[H]
\centering
\begin{tabular}{lll}
\toprule
\textbf{Student} & \textbf{Authored (O+C)} & \textbf{Assigned (C only)}\\
\midrule
Ethan Walsh & Num & Num \\
Colin Chambachan & Num & Num \\
Matthew Nesbitt & Num & Num \\
Jinwoo Hong & Num & Num \\
Gary Qin & Num & Num \\
\bottomrule
\end{tabular}
\end{table}

\wss{If needed, an explanation for the counts can be provided here.}

\section{CICD via GitHub Actions}

\wss{Say how CICD technology is used and will be used in your project.}

\wss{Provide links to your CICD yaml files}

\section{Custom Documents (CDs)}

The Custom Documents are as follows:

\begin{itemize}
    \item \textbf{Fall: Design Thinking Report} - This report was used to inform the team of user interface related decisions and to get feedback from peers. It will not be changed.
    \item \textbf{Winter: Performance Report} - This report will be used to gain more insight into the performance of the system as well as potential improvements. At this time, the team is in agreement
    that it should not be changed to another option.
\end{itemize}

\section{Team Charter Trigger Items}

\wss{Provide a summary of the quantified triggers identified in the team's
charter.}

\wss{Provide a list of any violations of the triggers.  If the team wishes, the
violations can be summarized on aggregate, instead of naming specific team
members.}

\wss{Provide a plan to address the violations.  This could include revising the
triggers, if they are found to be too weak, strong or ambiguous.}

\section{Additional Productivity Metrics}

\wss{If your team has additional metrics of productivity, please feel free to
add them to this report.  If not, please explicitly state that there are no additional
metrics.}

\wss{Additional metrics can include things like code reviews done, pull requests
created, count of joining meetings late, count of number of times contributions
had to be corrected, number of internal deadlines missed, test cases written, etc.}

\wss{We are looking for data on these metrics, not just a list of additional
metrics the team is planning on tracking.  However, if all you have is a plan,
please share it here.}

\end{document}